\documentclass[11pt,a4paper]{article}
\usepackage{fontspec,amsmath,amssymb,amsthm,microtype}
\setmainfont{lmroman10-regular.otf}[BoldFont=lmroman10-bold.otf,ItalicFont=lmroman10-italic.otf,BoldItalicFont=lmroman10-bolditalic.otf]
\usepackage{xcolor}
\usepackage[most]{tcolorbox}
\usepackage{mathrsfs}
\providecommand{\lt}{<}
\providecommand{\gt}{>}
\setmonofont{DejaVuSansMono.ttf}[BoldFont=DejaVuSansMono-Bold.ttf,ItalicFont=DejaVuSansMono-Oblique.ttf,BoldItalicFont=DejaVuSansMono-BoldOblique.ttf,Scale=0.86]
\newfontfamily\leanfallback{STIXTwoMath-Regular.otf}[Scale=0.86]
\definecolor{leanrule}{gray}{0.62}
\definecolor{leanback}{gray}{0.962}
\definecolor{leantitle}{gray}{0.22}
\definecolor{leanmuted}{gray}{0.38}
\newtcolorbox{leanbox}[1]{enhanced,breakable,colback=leanback,colframe=leanrule,
  boxrule=0.4pt,arc=0.6mm,left=2.4mm,right=2.4mm,top=1.8mm,bottom=1.8mm,
  before skip=0.8em,after skip=1em,fonttitle=\small\bfseries,coltitle=leantitle,
  colbacktitle=leanback,titlerule=0.3pt,title={#1}}
\newlength{\leancharwidth}
\newenvironment{leancode}{\par\smallskip\begingroup\ttfamily\footnotesize\raggedright
  \settowidth{\leancharwidth}{M}\setlength{\parskip}{0pt}\setlength{\parindent}{0pt}
  \frenchspacing}{\par\endgroup\smallskip}
\newcommand{\leanline}[2]{\hangindent=\dimexpr#1\leancharwidth+4\leancharwidth\relax
  \hangafter=1\noindent\hspace*{#1\leancharwidth}#2\par}
\newcommand{\leansource}[1]{\par{\footnotesize\color{leanmuted}#1}\par}
\newcommand{\leanlabel}[1]{\par\smallskip\noindent{\small\bfseries #1}\par}
\usepackage[margin=25mm]{geometry}
\usepackage[unicode,hidelinks]{hyperref}
\hypersetup{pdfauthor={LangevinLean Project}}
\setlength{\emergencystretch}{3em}
\setlength{\parskip}{0.4em}
\setlength{\parindent}{0pt}
\allowdisplaybreaks[2]
\raggedbottom
\title{\texorpdfstring{The infimum of squared coupling estimates}{The infimum of squared coupling estimates}}
\author{LangevinLean Project}
\date{}
\begin{document}\maketitle
Draft mathematical exposition; these notes may contain mistakes.


Assume \(W_2(\mu,\nu)<\infty\) and \(q,c\ge0\). Suppose \(W_2(\mu',\nu')^2\le q\operatorname{cost}(p)^2+c\) holds for every finite-cost coupling \(p\) of \(\mu,\nu\). The same coefficient and forcing then bound the output squared Wasserstein distance in terms of the input squared distance. The proof below treats zero coefficients separately and otherwise uses the greatest-lower-bound property of the input cost infimum.



Earlier results used here: \cite[Section~3.18]{LangevinLean2026ULASC}; \cite[Section~3.33]{LangevinLean2026ULASC}.

\section{\texorpdfstring{The greatest-lower-bound argument}{The greatest-lower-bound argument}}


If \(q=0\), choose one finite-cost input coupling. If the output square is at most \(c\), nonnegativity of the other term already proves the claim. Otherwise the coefficient is positive and one may set \(r=\sqrt{(W_2(\mu',\nu')^2-c)/q}\). Each finite coupling estimate implies \(r\le\operatorname{cost}(p)\), and infinite-cost couplings satisfy the same extended lower bound. Hence \(r\le\inf_p\operatorname{cost}(p)=W_2(\mu,\nu)\). Squaring and multiplying by \(q\) gives

\[W_2(\mu',\nu')^2\le qW_2(\mu,\nu)^2+c.\]

In the notation \(A=W_2(\mu',\nu')^2\), the hypothesis is \(A\le qC^2+c\) for all finite costs. The case \(q=0\) gives \(A\le c\) directly. The case \(A\le c\) follows using \(q\ge0\). Otherwise \(r=\sqrt{(A-c)/q}\) is real and positive, and all coupling costs are at least \(r\).



Equivalently, costs may approach the finite infimum from above; continuity of squaring makes their squared costs tend to \(W_2(\mu,\nu)^2\). The fixed output bound survives that limit. These arguments require neither an optimal coupling nor division by a coefficient without first proving its positivity. One finite-cost comparison also establishes output finiteness.

\begin{leanbox}{Lean: \texttt{wasserstein2\_\allowbreak{}sq\_\allowbreak{}bound\_\allowbreak{}of\_\allowbreak{}finite\_\allowbreak{}couplings}}
\leansource{\texttt{wasserstein2\_\allowbreak{}sq\_\allowbreak{}bound\_\allowbreak{}of\_\allowbreak{}finite\_\allowbreak{}couplings} \quad ProximalLangevin/\allowbreak{}SplittingTransport.lean, lines 21–58}
\begin{leancode}
\leanline{0}{theorem\ wasserstein2\_sq\_bound\_of\_finite\_couplings}
\leanline{4}{(μ\ ν\ μ'\ ν'\ :\ Measure\ E)\ (q\ c\ :\ ℝ)\ (hq\ :\ 0\ ≤\ q)\ (\_hc\ :\ 0\ ≤\ c)}
\leanline{4}{(hfinite\ :\ ULA.wasserstein2\ μ\ ν\ ≠\ ⊤)}
\leanline{4}{(hfamily\ :\ ∀\ p\ :\ ULA.Coupling\ μ\ ν,\ ULA.couplingCost\ p.joint\ ≠\ ⊤\ →}
\leanline{6}{ULA.wasserstein2Real\ μ'\ ν'\ \textasciicircum{}\ 2\ ≤\ q\ *\ (ULA.couplingCost\ p.joint).toReal\ \textasciicircum{}\ 2\ +\ c)\ :}
\leanline{4}{ULA.wasserstein2Real\ μ'\ ν'\ \textasciicircum{}\ 2\ ≤\ q\ *\ ULA.wasserstein2Real\ μ\ ν\ \textasciicircum{}\ 2\ +\ c}
\end{leancode}
\leanlabel{How the Lean proof follows the argument}
The operation \texttt{le\_\allowbreak{}i\allowbreak{}Inf} now proves \(r\le W_2(\mu,\nu)\) directly from the definition as an infimum. The finite-distance hypothesis justifies passing to real values. Squaring and multiplying by \(q\) gives \(A\le qW_2(\mu,\nu)^2+c\). This is why the proof does not need existence of an optimal transport plan.



Formal correspondence: wasserstein2\_sq\_bound\_of\_finite\_couplings : A uniform squared estimate for actual finite-cost couplings passes to the square of the genuine Wasserstein distance.



\texttt{wasserstein2\_\allowbreak{}sq\_\allowbreak{}bound\_\allowbreak{}of\_\allowbreak{}finite\_\allowbreak{}couplings}: the cases are separated by \texttt{by\_\allowbreak{}cases hq0 : q = 0}; the intermediate facts include \texttt{hqpos}, \texttt{hnum}, \texttt{r}, \texttt{hs}, \texttt{hreal}, in source order; the source steps apply or rewrite \texttt{zero\_\allowbreak{}mul}, \texttt{zero\_\allowbreak{}add}, \texttt{trans\_\allowbreak{}eq}, \texttt{ENNReal.\allowbreak{}of\allowbreak{}Real\_\allowbreak{}to\allowbreak{}Real}, \texttt{ENNReal.\allowbreak{}to\allowbreak{}Real\_\allowbreak{}of\allowbreak{}Real}. The remaining scalar rearrangements are checked by \texttt{field\_\allowbreak{}simp}, \texttt{linarith}, \texttt{nlinarith}.
\end{leanbox}


Earlier results used here: \cite[Section~3.17.1]{LangevinLean2026ULASC}; \cite[Section~3.17.2]{LangevinLean2026ULASC}; \cite[Section~3.17.3]{LangevinLean2026ULASC}; \cite[Section~3.17.4]{LangevinLean2026ULASC}; \cite[Section~3.18.2]{LangevinLean2026ULASC}.

\begin{thebibliography}{ULA-LSI}
\bibitem[ULA-SC]{LangevinLean2026ULASC}
LangevinLean Project. \emph{Unadjusted Langevin under Strong Convexity}. LL-01, 2026. Draft report, content revision d50c0f64815b. \url{https://odakyildiz.github.io/LangevinLean/docs/pdf/UnadjustedLangevin/FullProof.pdf}.
\end{thebibliography}

\end{document}
